% Generated draft: requires mathematical and source review.
\documentclass[11pt]{article}
\usepackage[margin=25mm]{geometry}
\usepackage{fontspec}
\setmainfont{Latin Modern Roman}
\usepackage{amsmath,amssymb,mathtools,unicode-math}
\setmathfont{Latin Modern Math}
\usepackage{enumitem,needspace,xurl,hyperref}
\setlength{\parindent}{0pt}
\setlength{\parskip}{5pt}
\setlength{\emergencystretch}{3em}
\newcommand{\N}{\mathbb{N}}
\newcommand{\Z}{\mathbb{Z}}
\newcommand{\Q}{\mathbb{Q}}
\newcommand{\R}{\mathbb{R}}
\newcommand{\C}{\mathbb{C}}
\newcommand{\F}{\mathbb{F}}
\newcommand{\A}{\mathbb{A}}
\newcommand{\PP}{\mathbb{P}}
\newcommand{\E}{\mathbb{E}}
\newcommand{\eps}{\varepsilon}
\newcommand{\dd}{\,\mathrm{d}}
\begin{document}
\section*{UnsolvedMath ID: 701365}
\textbf{Source problem code:} catalogue entry
\subsection{Definitions and mathematical statement}
Fix an odd prime \(p\) and a positive integer \(s\). Write \(\Z_p=\varprojlim_n\Z/p^n\Z\) for the ring of \(p\)-adic integers. For a finite \(p\)-group \(P\), define its coclass by
\[
\operatorname{cc}(P)=\log_p|P|-\operatorname{cl}(P),
\]
where \(\operatorname{cl}(P)\) is its nilpotency class, with the trivial group assigned class zero.

The supplied record does not define a \(p\)-adic space group. For a conditional precise formulation, use the interpretation of a pro-\(p\) group \(G\) admitting an exact sequence of topological groups
\[
1\longrightarrow T\longrightarrow G\longrightarrow P\longrightarrow 1,
\]
where \(T\cong(\Z_p^d,+)\) for some finite integer \(d\geq1\), \(P\) is a finite \(p\)-group, and the induced conjugation action \(P\longrightarrow\operatorname{Aut}_{\Z_p}(T)\) is faithful. Here \(T\) is the translation subgroup and \(P\) is the point group. No splitting of the extension is assumed. This interpretation requires confirmation against the cited source.

For a pro-\(p\) group \(G\), put \(\gamma_1(G)=G\) and \(\gamma_{n+1}(G)=\overline{[\gamma_n(G),G]}\). Use the convention that \(G\) has finite coclass if, for all sufficiently large \(n\), the quotient \(G/\gamma_n(G)\) is finite and the integers \(\operatorname{cc}(G/\gamma_n(G))\) are eventually constant at a finite value.

With these conventions, the proposed assertion is
\[
\forall\,p\text{ odd prime}\quad\forall\,s\in\Z_{\geq1},\qquad
\#\left\{[G]_{\mathrm{top}}:\begin{array}{l}
G\text{ is a }p\text{-adic space group of finite coclass},\\
\operatorname{cc}(P)=s\text{ for its point group }P
\end{array}\right\}<\infty,
\]
where \([G]_{\mathrm{top}}\) denotes the continuous group isomorphism class of \(G\). The source does not explicitly specify this equivalence relation; it must also be confirmed. The finite coclass of \(G\) is allowed to vary: the assertion imposes no fixed value or additional bound on it.
\subsection{Short English statement}
For every fixed odd prime \(p\) and positive integer \(s\), prove that only finitely many \(p\)-adic space groups of finite coclass have point group of coclass \(s\), interpreted here up to continuous group isomorphism. The precise space-group convention and counting equivalence require source confirmation.
\subsection{Sources}
\begin{itemize}
\item Local UnsolvedMath registry: permanent ID \(701365\), exact source problem code \texttt{catalogue entry}, title \emph{Kourovka Notebook Problem 17.63}; supplied upstream revision \texttt{d371ee917ac9647fbb06cd6fc529deec6f3dda82}.
\item Attributed in the supplied record to C. R. Leedham-Green, proposed in 2010; original problem number 17.63 in \emph{The Kourovka Notebook: Unsolved Problems in Group Theory}, 17th issue (2010). The record also identifies the 21st issue edition (2026), \texttt{arXiv:1401.0300v48}. Its collection description names E. I. Khukhro and V. D. Mazurov as editors.
\item The registry labels the problem \texttt{open}. This is supplied source metadata, not an independently verified assessment. No separate literature-verification date or literature-source links are supplied, and the cited editions have not been independently consulted for this draft.
\end{itemize}
\end{document}
